% ===================================================================
% data/kennzahlen.tex
% EINZIGE Quelle aller Zahlenwerte des Dokuments.
% Rohwerte im englischen Format (Punkt als Dezimaltrenner) -
% die Formatierung mit Dezimalkomma erfolgt in preamble.tex.
% Jedes Makro traegt seine Quelle und Seitenzahl im Kommentar.
%
% Alle Betraege in Mio. EUR, sofern nicht anders vermerkt. Die Werte sind
% in voller im Bericht veroeffentlichter Genauigkeit hinterlegt (Aumann
% weist in TEUR aus); \MioEUR rundet fuer den Fliesstext auf eine
% Nachkommastelle, \MioEURexakt gibt den vollen Wert im Rechenweg aus.
% ===================================================================

% --- Aufgabe 2: Jahresabschluss 2024 -------------------------------
% Konzernkennzahlen laut Uebersicht "Aumann in figures"
\newcommand{\UmsatzZFIII}{289.606}            % GB 2024, S. 2 (289.606 TEUR)
\newcommand{\UmsatzZFIV}{312.346}             % GB 2024, S. 2 (312.346 TEUR)
\newcommand{\AuftragseingangZFIII}{339.379}   % GB 2024, S. 2 (339.379 TEUR)
\newcommand{\AuftragseingangZFIV}{200.057}    % GB 2024, S. 2 (200.057 TEUR)

% Aumann weist EBITDA und bereinigtes (adjusted) EBITDA getrennt aus.
% Als "operatives EBITDA" im Sinne der Aufgabenstellung wird das bereinigte
% EBITDA gefuehrt; das berichtete EBITDA steht zur Einordnung daneben.
\newcommand{\EbitdaBerichtetZFIII}{20.647}    % GB 2024, S. 2 (EBITDA)
\newcommand{\EbitdaBerichtetZFIV}{35.804}     % GB 2024, S. 2 (EBITDA)
\newcommand{\EbitdaMargeBerichtetZFIII}{7.1}  % GB 2024, S. 2 (EBITDA margin)
\newcommand{\EbitdaMargeBerichtetZFIV}{11.5}  % GB 2024, S. 2 (EBITDA margin)
\newcommand{\EbitdaOpZFIII}{21.294}           % GB 2024, S. 2 (Adj. EBITDA)
\newcommand{\EbitdaOpZFIV}{36.417}            % GB 2024, S. 2 (Adj. EBITDA)
\newcommand{\EbitdaMargeZFIII}{7.4}           % GB 2024, S. 2 (Adj. EBITDA margin)
\newcommand{\EbitdaMargeZFIV}{11.7}           % GB 2024, S. 2 (Adj. EBITDA margin)
\newcommand{\BereinigungsbetragZFIV}{0.613}   % berechnet: EbitdaOpZFIV - EbitdaBerichtetZFIV

\newcommand{\EbitZFIV}{29.454}                % GB 2024, S. 2
\newcommand{\NettogewinnZFIII}{9.583}         % GB 2024, S. 2 (Consolidated net profit)
\newcommand{\NettogewinnZFIV}{21.506}         % GB 2024, S. 2 (Consolidated net profit)

\newcommand{\DividendeJeAktieZFIV}{0.22}      % GB 2024, S. 8 (Gewinnverwendungsvorschlag)
\newcommand{\KursSilvesterZFIV}{10.62}        % GB 2024, S. 13 (XETRA-Schlusskurs 30.12.2024)
\newcommand{\DividendenrenditeZFIV}{2.07}     % berechnet: Dividende / Schlusskurs
\newcommand{\NettoliquiditaetZFIV}{138.2}     % GB 2024, S. 5

% --- Basiswerte 2025 (Bezugsbasis fuer Aufgabe 6) ------------------
\newcommand{\UmsatzZFV}{203.985}              % GB 2025, S. 2
\newcommand{\AuftragseingangZFV}{147.486}     % GB 2025, S. 2
\newcommand{\EbitdaBerichtetZFV}{28.215}      % GB 2025, S. 2 (EBITDA)
\newcommand{\EbitdaOpZFV}{27.278}             % GB 2025, S. 2 (Adj. EBITDA)
\newcommand{\EbitdaMargeZFV}{13.4}            % GB 2025, S. 2 (Adj. EBITDA margin)
\newcommand{\NettoliquiditaetZFV}{148.1}      % GB 2025, S. 4
\newcommand{\KursSilvesterZFV}{12.32}         % GB 2025, S. 13 (XETRA-Schlusskurs 2025)

% --- Stammdaten der Aktie ------------------------------------------
% Quelle der Web-Angaben: aumann.com/en/investor-relations/shares,
% abgerufen am 27.07.2026.
\newcommand{\ISIN}{DE000A2DAM03}              % IR-Webseite "Shares"
\newcommand{\WKN}{A2DAM0}                     % IR-Webseite "Shares"
\newcommand{\AktienGesamt}{12917048}          % IR-Webseite "Shares"
\newcommand{\AktienGesamtAlt}{15250000}       % GB 2025, S. 52 (vor Kapitalherabsetzungen)
\newcommand{\AktienGesamtMittel}{14345231}    % GB 2025, S. 52 (nach erster Herabsetzung)
% Die Angaben der IR-Webseite tragen ungleiche Stichtage: der MBB-Anteil ist
% mit "* as of 31 December 2025" datiert, der Streubesitz traegt kein Datum.
% Der frueher hier gefuehrte Stand (44,49 % / 45,5 %) stammt aus der Fassung
% der Seite vor der Kapitalherabsetzung und bezog sich noch auf
% \AktienGesamtMittel Aktien; Beleg: Internet Archive, Fassung vom 14.09.2025.
\newcommand{\AnteilMBB}{47.81}                % IR-Webseite "Shares", Stand 31.12.2025 (in %)
\newcommand{\FreefloatAnteil}{42.9}           % IR-Webseite "Shares", ohne Stichtag (in %)
\newcommand{\RestanteilUngeklaert}{9.29}      % berechnet: 100 - AnteilMBB - FreefloatAnteil
\newcommand{\AnteilMBBAlt}{44.49}             % IR-Webseite "Shares", archivierte Fassung 14.09.2025 (in %)
\newcommand{\FreefloatAnteilAlt}{45.5}        % ebenda (in %)

% Eigene Aktien aus dem Rueckkaufangebot 2026 (Groessenordnung der Differenz)
\newcommand{\EigeneAktienStueck}{1291200}     % Aumann AG, Meldung nach Paragraf 40 Abs. 1 Satz 2 WpHG vom 16.07.2026
\newcommand{\EigeneAktienAnteil}{9.996}       % ebenda (in % der Stimmrechte)
\newcommand{\RueckkaufPreisZFVI}{17.80}       % IR-Webseite "Share buyback" (Angebotspreis 2026, angehoben)
\newcommand{\IPOKurs}{42.00}                  % IR-Webseite "Shares" (Emissionspreis 2017)
\newcommand{\MarktkapZFV}{159.138}            % berechnet: AktienGesamt * KursSilvesterZFV
\newcommand{\MitarbeiterZFIV}{891}            % GB 2025, S. 93
\newcommand{\MitarbeiterZFV}{773}             % GB 2025, S. 93

% --- Aufgabe 2.7: XETRA-Jahresschlusskurse -------------------------
% Der Schlusskurs 2022 ist im GB 2023 als Vorjahreswert ausgewiesen.
% Die Ketten der Vorjahresangaben in GB 2023, 2024 und 2025 sind konsistent.
% Der Wert 2021 stammt allein aus Yahoo Finance: Der GB 2022 nennt keinen
% Jahresschlusskurs (geprueft am 29.07.2026). Er ist damit als einziger der
% fuenf nicht gegen eine Primaerquelle gepruefet; im Dokument ist das benannt.
\newcommand{\KursSilvesterZFI}{13.68}         % Yahoo Finance (ohne Primaerquellen-Gegenprobe)
\newcommand{\KursSilvesterZFII}{11.48}        % GB 2023, S. 14 (Vorjahresschlusskurs)
\newcommand{\KursSilvesterZFIII}{18.58}       % GB 2023, S. 14
\newcommand{\KursDeltaZFII}{-16.1}            % berechnet: (11.48-13.68)/13.68; erbt den Vorbehalt zu 2021
\newcommand{\KursDeltaZFIII}{61.8}            % GB 2023, S. 14 (Veraenderung 2023)
\newcommand{\KursDeltaZFV}{16.0}              % GB 2025, S. 13 (Veraenderung 2025)

% --- Jahreshoechst- und Jahrestiefstkurse (Intraday) ---------------
% Quelle: Yahoo Finance, Symbol AAG.DE (Handelsplatz XETRA), Monatsbars,
% abgerufen am 29.07.2026. Herleitung und Vorbehalte: data/aktienkurs.csv,
% dieselbe Datei speist die Abbildung. Es sind Intraday-Extrema, nicht
% schlusskursbasierte Extrema.
% 2021 dient allein als Ausgangswert der Abbildung; seine Extrema liegen vor
% dem Betrachtungszeitraum und werden daher nicht gefuehrt.
\newcommand{\KursHochZFII}{17.68}             % 2022
\newcommand{\KursTiefZFII}{10.10}             % 2022
\newcommand{\KursHochZFIII}{18.98}            % 2023
\newcommand{\KursTiefZFIII}{11.28}            % 2023
\newcommand{\KursHochZFIV}{18.96}             % 2024
\newcommand{\KursTiefZFIV}{9.42}              % 2024
\newcommand{\KursHochZFV}{14.66}              % 2025
\newcommand{\KursTiefZFV}{9.87}               % 2025

% Handelstage der Extrema. Sie stehen in Tabelle tab:kursextrema, weil die
% Abbildung sie nicht aufloesen kann: Hoch 2023 (28.12.), Schluss 2023
% (29.12.) und Hoch 2024 (02.01.) liegen fuenf Tage auseinander und fallen
% auf einer Vierjahresachse zu einer Spitze zusammen.
\newcommand{\DatumHochZFII}{23.03.2022}
\newcommand{\DatumTiefZFII}{13.10.2022}
\newcommand{\DatumHochZFIII}{28.12.2023}
\newcommand{\DatumTiefZFIII}{05.01.2023}
\newcommand{\DatumHochZFIV}{02.01.2024}
\newcommand{\DatumTiefZFIV}{19.11.2024}
\newcommand{\DatumHochZFV}{12.05.2025}
\newcommand{\DatumTiefZFV}{04.03.2025}
\newcommand{\KursSpanneZFIV}{50.3}            % berechnet: 1 - KursTiefZFIV/KursHochZFIV (Verlust vom Hoch zum Tief 2024, in %)

% --- Aufgabe 3: Segmente, Strategie und Prognose -------------------
\newcommand{\EigenkapitalZFIV}{201.7}         % GB 2024, S. 5
\newcommand{\AuftragsbestandZFIV}{184.0}      % GB 2024, S. 5
\newcommand{\AuftragsbestandDeltaZFIV}{-39.3} % GB 2024, S. 5
\newcommand{\AnteilEMobility}{80}             % GB 2024, S. 10 (Anteil am Geschaeft, in %)
\newcommand{\StandorteGesamt}{6}              % GB 2024, S. 9 ("six locations")
\newcommand{\AnteilUSGesellschaft}{100}       % GB 2024, S. 52 (Beteiligungsquote in %)

% Segment E-Mobility
\newcommand{\UmsatzEMobilityZFIV}{258.530}    % GB 2024, S. 2 ("thereof E-mobility")
\newcommand{\UmsatzEMobilityDeltaZFIV}{12.8}  % GB 2024, S. 13
\newcommand{\EbitdaEMobilityZFIII}{17.1}      % GB 2024, S. 13
\newcommand{\EbitdaEMobilityZFIV}{33.8}       % GB 2024, S. 13
\newcommand{\EbitdaMargeEMobilityZFIV}{13.1}  % GB 2024, S. 13

% Segment Next Automation (vormals Classic)
\newcommand{\UmsatzNextAutoZFIII}{60.5}       % GB 2024, S. 13
\newcommand{\UmsatzNextAutoZFIV}{53.8}        % GB 2024, S. 13
\newcommand{\UmsatzNextAutoZFV}{40.2}         % GB 2025, S. 13
\newcommand{\EbitdaNextAutoZFIII}{6.2}        % GB 2024, S. 13
\newcommand{\EbitdaNextAutoZFIV}{5.8}         % GB 2024, S. 13
\newcommand{\EbitdaNextAutoZFV}{5.1}          % GB 2025, S. 13
\newcommand{\EbitdaMargeNextAutoZFIII}{10.2}  % GB 2024, S. 13
\newcommand{\EbitdaMargeNextAutoZFIV}{10.8}   % GB 2024, S. 13
\newcommand{\EbitdaMargeNextAutoZFV}{12.8}    % GB 2025, S. 13
\newcommand{\AuftragseingangNextAutoZFIV}{36.6} % GB 2024, S. 13
\newcommand{\AuftragseingangNextAutoZFV}{56.5}  % GB 2025, S. 13

% Kapitalallokation
\newcommand{\RueckkaufAktienMax}{1434523}     % GB 2024, S. 5 (oeffentliches Angebot)
\newcommand{\RueckkaufPreis}{12.37}           % GB 2024, S. 5 (Angebotspreis je Aktie)

% Prognose 2025 aus dem Geschaeftsbericht 2024
\newcommand{\PrognoseUmsatzZFVUnten}{210}     % GB 2024, S. 5
\newcommand{\PrognoseUmsatzZFVOben}{230}      % GB 2024, S. 5
\newcommand{\PrognoseMargeZFVUnten}{8}        % GB 2024, S. 5
\newcommand{\PrognoseMargeZFVOben}{10}        % GB 2024, S. 5
\newcommand{\EbitdaMargeBerichtetZFV}{13.8}   % GB 2025, S. 2 (berichtete EBITDA-Marge)
\newcommand{\PrognoseAbweichungUmsatz}{-2.9}  % berechnet: (203.985-210)/210

% --- Abgeleitete Veraenderungsraten (in Prozent) -------------------
% Berechnet aus den exakten Werten oben; entspricht der Spalte
% "Delta 2024 / 2023" des Geschaeftsberichts.
\newcommand{\UmsatzDeltaZFIV}{7.9}            % (312.346-289.606)/289.606
\newcommand{\AuftragseingangDeltaZFIV}{-41.1} % (200.057-339.379)/339.379
\newcommand{\NettogewinnDeltaZFIV}{124.4}     % (21.506-9.583)/9.583
\newcommand{\EbitdaOpDeltaZFIV}{71.0}         % GB 2024, S. 2 (Delta-Spalte Adj. EBITDA)
\newcommand{\KursDeltaZFIV}{-42.9}            % GB 2024, S. 13 (Kursrueckgang 2024)
\newcommand{\UmsatzDeltaZFV}{-34.7}           % (203.985-312.346)/312.346

% --- Aufgabe 1.5: dokumentierte Analysten-Ratingaktionen -----------
% Einzelne, datierte Ratingaktionen (nicht der Konsens-Aggregatoren,
% die am 02.08.2026 keine Daten zeigten - siehe quelle:boerse-express /
% quelle:boersede-empfehlungen). Sekundaerquellen (Finanzportale, nicht
% Aumann selbst); Schlagzeilen/Suchtreffer mehrfach uebereinstimmend
% geprueft, kein Volltext-Fetch moeglich (403), ausser EQUI.TS.
\newcommand{\KurszielBerenbergVorHochstufung}{18} % vor Aug. 2023, GB/Presse-Schlagzeile
\newcommand{\KurszielBerenbergHochstufung}{21}    % nach Hochstufung Halten->Kaufen, Aug. 2023
\newcommand{\KurszielBerenbergAnhebung}{22}       % Anhebung innerhalb Kaufen, Nov. 2023
\newcommand{\KurszielBerenbergRueckstufung}{17}   % Rueckstufung Kaufen->Halten, Jul. 2024

% Analysten-Konsens (stockanalysis.com, 2 Analysten, Stand 02.08.2026)
\newcommand{\KurszielKonsensDurchschnitt}{14.85}
\newcommand{\KurszielKonsensTief}{13.70}
\newcommand{\KurszielKonsensHoch}{16.00}

% --- Anker fuer die eigene Schaetzung Teil 5 (F&E-Intensitaet) -----
% Aktivierte Entwicklungskosten und ihr Anteil am Umsatz, wie im
% Bericht ausgewiesen (nicht aus dem Umsatz zurueckgerechnet, siehe
% docs/befunde-und-entscheidungen.md zum Umgang mit gedruckten Werten).
\newcommand{\EntwicklungskostenAktiviertZFIII}{2.7} % GB 2023, S. 14
\newcommand{\FEQuoteZFIII}{1.5}               % GB 2023, S. 14 (Anteil am Umsatz, in %)
\newcommand{\EntwicklungskostenAktiviertZFV}{1.9}   % GB 2025, S. 13
\newcommand{\FEQuoteZFV}{0.8}                 % GB 2025, S. 13 (Anteil am Umsatz, in %)

% --- Aufgabe 6: Parameter der Margen-Sensitivitaet -----------------
% Die Marge Ende 2028 haengt allein davon ab, welcher Umsatz im Nenner
% steht. Der Basisfall schreibt den Umsatz 2025 konstant fort - und der
% war gegenueber 2024 um \UmsatzDeltaZFV Prozent eingebrochen. Die
% Szenariorechnung steht in data/marge_sensitivitaet.tex.
\newcommand{\MargenWachstumSzenario}{3}       % Annahme: Umsatzwachstum in % p.a.
\newcommand{\MargenSzenarioJahre}{3}          % Wachstumsjahre 2026-2028
